\documentclass[11pt,a4paper]{article}

\usepackage[T1]{fontenc}
\usepackage[utf8]{inputenc}
\usepackage[margin=1in]{geometry}
\usepackage{amsmath, amssymb}
\usepackage{graphicx}
\usepackage{booktabs}
\usepackage{hyperref}
\usepackage{xcolor}
\usepackage{subcaption}
\usepackage{float}
\usepackage{cite}
\usepackage{microtype}
\usepackage{parskip}
\usepackage{titlesec}
\usepackage{fancyhdr}
\usepackage{enumitem}
\usepackage{abstract}
\usepackage{iftex}
\ifXeTeX
  \usepackage{fontspec}
  \setmainfont{FreeSerif}
  \newfontfamily\devfont{FreeSerif}
  \newcommand{\dev}[1]{{\devfont #1}}
\else\ifLuaTeX
  \usepackage{fontspec}
  \setmainfont{FreeSerif}
  \newfontfamily\devfont{FreeSerif}
  \newcommand{\dev}[1]{{\devfont #1}}
\else
  \newcommand{\dev}[1]{#1}
\fi\fi

\hypersetup{
    colorlinks=true,
    linkcolor=blue!60!black,
    citecolor=blue!60!black,
    urlcolor=blue!60!black
}

\titleformat{\section}{\large\bfseries}{{\thesection.}}{0.5em}{}
\titleformat{\subsection}{\normalsize\bfseries}{{\thesubsection}}{0.5em}{}

\pagestyle{fancy}
\fancyhf{}
\rhead{\small SpeakPay — Technical Report}
\lhead{\small \thepage}
\renewcommand{\headrulewidth}{0.4pt}

% -- Title -------------------------------------------------------------------
\title{
    \vspace{-1cm}
    {\LARGE \textbf{SpeakPay: Domain-Adaptive LoRA Fine-Tuning of Whisper\\[4pt]
    for Low-Resource Nepali Financial Speech Recognition}}
}
\author{
    Biraj Subedi \\
    \small Advanced College of Engineering and Management, Tribhuvan University \\
    \small \href{https://github.com/subedibiraj/speakpay}{github.com/subedibiraj/speakpay} \\
    \small \texttt{contact@biraj-subedi.com.np}
}
\date{Technical Report $\cdot$ \today}

% ─────────────────────────────────────────────────────────────────────────────
\begin{document}
\maketitle

% -- Abstract ----------------------------------------------------------------
\begin{abstract}
Mobile payment applications in Nepal are graphically mediated and largely
inaccessible to visually impaired users. This report presents SpeakPay, a
voice-first digital wallet, and documents the central technical
contribution: a controlled study of domain adaptation for low-resource
financial speech recognition. We introduce NepFinSpeech-403, a 403-utterance
dataset of Nepali financial voice commands (send, load, and balance
operations spanning 237 unique numerals), and fine-tune Whisper large-v2
test set, the domain-adapted model reduces Word Error Rate from 129.95\%
(zero-shot baseline) to 42.58\% --- a 67.2\% relative reduction --- and
improves Devanagari numeral recognition accuracy from 0.0\% to 73.9\%
averaged across utterances that contain numerals. Crucially, we find that
word-level metrics severely understate the practical task-level impact:
domain adaptation improves the Transaction Success Rate (Interpretation Error Rate)
from 1.67\% under the zero-shot model to 33.33\% under the adapted model. The
improvement is consistent at the individual-utterance level: the domain-adapted model
outperforms the zero-shot baseline on 59 of 60 test utterances (sign test,
$p < 10^{-17}$). A per-intent breakdown shows the gain holds across
all three command types (send, load, and a residual ``other'' category),
though the ``other'' category reveals a labeling gap in our intent
taxonomy --- discussed in Section~6. The trained system, including a
self-hosted inference API and a voice-first web application, is deployed
and publicly accessible. All code, the dataset, model weights, and this
report are released openly at
\url{https://github.com/subedibiraj/speakpay}.
\end{abstract}

\vspace{0.5em}
\hrule
\vspace{1em}

% ── 1. Introduction ──────────────────────────────────────────────────────────
\section{Introduction}

Mobile wallets (eWallets) such as eSewa and Khalti have become the dominant
mode of digital payment in Nepal. Their interfaces, however, are built
entirely around visual interaction --- icon-based navigation, on-screen
forms, and graphical confirmation flows. For Nepal's visually impaired
population (estimated at roughly 95{,}000 individuals~\cite{who2019}), this
makes independent use of digital finance effectively impossible without
sighted assistance.

Voice interaction is the natural accessibility solution, but two gaps stand
in the way. First, existing Nepali Automatic Speech Recognition (ASR)
systems are trained and evaluated on general speech~\cite{regmi2021,
paudel2023}, not financial commands, which carry domain-specific demands:
dense numeral sequences (transfer amounts, account balances), proper nouns
(bank names, recipient names), and a narrow set of recurring sentence
structures. Second, no public Nepali financial speech dataset existed prior
to this work, making it impossible to even measure how badly general-purpose
models perform on this domain, let alone improve on it.

This report documents the construction of such a dataset and a controlled
comparison of zero-shot and domain-adapted ASR on it. The project began as
an undergraduate group project with no rigorous evaluation, no public
dataset, and no reproducible benchmark. This version is a complete
independent rewrite: a newly assembled and characterized dataset, a
reproducible LoRA fine-tuning pipeline, a statistically grounded evaluation,
and a deployed application.

% ── 2. Related Work ──────────────────────────────────────────────────────────
\section{Related Work}

\textbf{Nepali ASR.} Regmi and Bal~\cite{regmi2021} built an end-to-end
Nepali ASR system with ESPnet on OpenSLR data, reporting a Character Error
Rate (CER) of 10.3\% on 159k general-domain utterances. Paudel et
al.~\cite{paudel2023} used a CNN-Transformer architecture, achieving CER of
11.14\% on a similarly general corpus. Both systems are evaluated on
read speech or broad-domain utterances; neither reports performance on
financial or transactional language, which differs substantially in
vocabulary distribution and numeral density.

\textbf{Whisper and parameter-efficient fine-tuning.} Radford et
al.~\cite{radford2022} introduced Whisper, an encoder-decoder Transformer
trained on 680{,}000 hours of weakly supervised multilingual audio.
Whisper's zero-shot performance on Nepali is poor relative to
high-resource languages; we quantify this directly for the financial
domain in Section~5. Hu et al.~\cite{hu2022lora} introduced LoRA, which
freezes the base model and injects trainable low-rank decomposition
matrices into attention and feed-forward layers, reducing trainable
parameters by over 99\% relative to full fine-tuning while matching or
exceeding its quality on many tasks. This makes adaptation feasible on a
single consumer GPU and, critically, on a dataset far too small to safely
full-fine-tune a 1.55-billion-parameter model.

\textbf{Accessible FinTech.} Concurrent work on accessible voice payments \cite{accessible_mobile_money2026, voicebiometric_upi2026} has focused on system architecture, USSD automation, and biometric security for visually impaired users. We focus instead on the ASR domain-adaptation and evaluation-metric questions underlying the language layer these systems depend on, using Nepali --- a language with no existing accessible financial voice interface --- as a case study.

Our work does not propose a new architecture or training method. It
presents a controlled, statistically evaluated study of whether
domain-specific LoRA adaptation closes the gap between a general-purpose
ASR model and the specific requirements of financial voice interaction in a
low-resource language --- and ships the result as a working application.

% ── 3. Dataset: NepFinSpeech-403 ─────────────────────────────────────────────
\section{Dataset: NepFinSpeech-403}

\subsection{Collection}

Audio was collected from undergraduate students and staff at Advanced
College of Engineering and Management (Tribhuvan University) through a
purpose-built web platform
(\url{https://bolanepal.netlify.app}). Contributors
recorded spoken Nepali financial commands guided by written prompts,
covering three operation types: fund transfers (recipient, amount,
optionally a financial institution), wallet load/deposit commands, and
balance enquiries. Transcripts were manually verified against the recorded
audio. Notably, contributor identity was not recorded per utterance during
collection, meaning speaker overlap between train and test splits cannot
be entirely ruled out.

\subsection{Statistics}

\begin{table}[H]
\centering
\caption{NepFinSpeech-403 composition}
\label{tab:dataset}
\begin{tabular}{lrr}
\toprule
\textbf{Split} & \textbf{Samples} & \textbf{\%} \\
\midrule
Train & 303 & 75\% \\
Validation & 40 & 10\% \\
Test (held-out) & 60 & 15\% \\
\midrule
\textbf{Total} & \textbf{403} & 100\% \\
\bottomrule
\end{tabular}
\quad
\begin{tabular}{lr}
\toprule
\textbf{Intent (heuristic)} & \textbf{Count} \\
\midrule
Send & 193 (47.9\%) \\
Balance & 61 (15.1\%) \\
Load & 56 (13.9\%) \\
Other / unlabeled & 93 (23.1\%) \\
\midrule
Unique Devanagari numerals & 237 \\
\bottomrule
\end{tabular}
\end{table}

The ``other'' category is a heuristic residual class assigned by keyword
matching, not a manually curated label; Section~6 shows it captures a
genuine, previously unaccounted-for intent (third-person ``funds received''
statements), which we treat as a finding rather than noise.

% ── 4. Method ─────────────────────────────────────────────────────────────────
\section{Method}

\textbf{Base model.} Whisper large-v2~\cite{radford2022}, 1.55B parameters,
encoder-decoder Transformer operating on log-Mel spectrograms.

\textbf{Adaptation.} LoRA~\cite{hu2022lora} with rank $r=32$, scaling
$\alpha=64$, applied to all query, key, value, output, and feed-forward
projection matrices (\texttt{q\_proj}, \texttt{k\_proj}, \texttt{v\_proj},
\texttt{out\_proj}, \texttt{fc1}, \texttt{fc2}) across every encoder and
decoder layer. This yields approximately 8M trainable parameters, under
0.6\% of the base model, while the remaining 99.4\% stays frozen.

\textbf{Why LoRA, not full fine-tuning.} With only 303 training utterances,
full fine-tuning of a 1.55B-parameter model risks both overfitting and
catastrophic forgetting of the base model's general speech competence.
LoRA's parameter efficiency is not merely a compute convenience here --- it
is the difference between a feasible and an infeasible experiment at this
dataset scale.

\textbf{Training configuration.} fp16 precision, no quantization, single
consumer GPU (RTX 3060, 12GB), effective batch size 16 via gradient
accumulation, learning rate $1\times10^{-4}$. We train for a maximum of 300
steps over the 303-utterance training set, corresponding to roughly 15.8 epochs.
Full configuration in \texttt{training/scripts/config.py} of the accompanying repository.

% ── 5. Results ─────────────────────────────────────────────────────────────────
\section{Results}

\subsection{Aggregate benchmark}

All models are evaluated on the same 60-utterance held-out test split,
never seen during training (fixed seed for reproducibility).

\begin{table}[H]
\centering
\caption{Benchmark results on the NepFinSpeech-403 test set}
\label{tab:results}
\begin{tabular}{lccc}
\toprule
\textbf{Model} & \textbf{WER\% $\downarrow$} & \textbf{CER\% $\downarrow$} & \textbf{NumAcc\% $\uparrow$} \\
\midrule
Whisper large-v2 (zero-shot) & 129.95 & 92.32 & 0.0 \\
Whisper large-v3-turbo (general Nepali) & --- & --- & --- \\
\textbf{Whisper large-v2 + LoRA (ours)} & \textbf{42.58} & \textbf{16.95} & \textbf{73.9} \\
\bottomrule
\end{tabular}
\end{table}

\textit{The general-domain Nepali fine-tune baseline failed to load during
this benchmark run; see Section~7 (Limitations) for discussion. The
remaining comparison is between the zero-shot base model and our
domain-adapted model.}

The zero-shot WER exceeds 100\%, which indicates the model's output is, on
average, longer than the reference --- consistent with insertion-heavy
hallucination rather than simple mistranscription. Manual inspection of
zero-shot predictions confirms this: outputs frequently contain
phonetically plausible but semantically incoherent Devanagari, including
malformed numeral sequences and occasionally repeated phrase fragments.
Domain adaptation does not merely improve transcription accuracy; it
recovers basic output coherence in a domain where the base model fails
qualitatively, not just quantitatively.

\subsection{Per-utterance significance}

Aggregate WER can be dominated by a small number of catastrophic failures.
We therefore additionally report a paired, per-utterance comparison: for
each of the 60 test utterances, we compute WER under both the zero-shot
and domain-adapted models and compare directly.

\begin{table}[H]
\centering
\caption{Paired per-utterance comparison (domain-adapted vs.\ zero-shot)}
\label{tab:paired}
\begin{tabular}{lr}
\toprule
Utterances where domain-adapted model improves & 59 / 60 \\
Utterances where zero-shot is better & 0 / 60 \\
Tied & 1 / 60 \\
Mean per-utterance WER reduction & 87.7 percentage points \\
Sign test (two-sided, $n=59$ non-tied pairs) & $p = 3.5\times10^{-18}$ \\
\bottomrule
\end{tabular}
\end{table}

The zero-shot model yielded a mean per-utterance WER of 130.33\% (95\% bootstrap CI [115.15, 152.96]), whereas the domain-adapted model achieved 42.61\% (95\% bootstrap CI [36.76, 48.59]).

The improvement is not driven by a handful of outliers: the domain-adapted
model is better on essentially every individual utterance in the test set,
and the sign test rejects the null hypothesis of no systematic difference
at an extremely high confidence level. Given the small absolute sample
size ($n=60$), we report the non-parametric sign test rather than relying
assumptions and outlier influence.

\subsection{Task-level Interpretation Accuracy}

While Word Error Rate is the standard ASR benchmark, it treats all word errors equally. In a financial voice interface, a wrong filler word is a harmless transcription error, but a mistranscribed digit is a real-money transaction failure. To measure practical utility, we evaluate both models using Interpretation Error Rate (IRER)~\cite{fu2022multitask}, defined here as the utterance-level Transaction Success Rate: the percentage of test utterances where the predicted transcript, when parsed by a rule-based slot extraction pipeline, yields exactly the same intent, amount, and recipient as the ground-truth transcript.

\begin{table}[H]
\centering
\caption{Task-level slot extraction accuracy (higher is better)}
\label{tab:task_eval}
\begin{tabular}{lrr}
\toprule
\textbf{Metric} & \textbf{Zero-shot} & \textbf{Ours (LoRA)} \\
\midrule
Intent Accuracy & 73.33\% & 98.33\% \\
Amount Exact Match & 12.28\% & 54.39\% \\
Recipient Exact Match & 5.26\% & 42.11\% \\
\midrule
\textbf{Transaction Success Rate (1 - IRER)} & \textbf{1.67\%} & \textbf{33.33\%} \\
\bottomrule
\end{tabular}
\end{table}

The domain adaptation increases the Transaction Success Rate from 1.67\% (effectively unusable) to 33.33\%. While 33.33\% is still too low for a fully autonomous financial system without user confirmation steps, the relative improvement over the zero-shot baseline is nearly 1900\% --- significantly larger than the 67\% relative improvement seen in WER alone. This confirms that word-level metrics heavily understate the task-level impact of domain adaptation for slot-critical applications.

\subsection{Per-intent breakdown}

\begin{table}[H]
\centering
\caption{WER by command intent (heuristic classification)}
\label{tab:intent}
\begin{tabular}{lrrr}
\toprule
\textbf{Intent} & \textbf{N} & \textbf{Zero-shot WER\%} & \textbf{Ours WER\%} \\
\midrule
Send     & 30 & 143.6 & 45.8 \\
Load     &  8 & 124.1 & 34.8 \\
Other    & 22 & 114.5 & 41.1 \\
\bottomrule
\end{tabular}
\end{table}

The improvement holds consistently across all three groups, with the
largest absolute gain on the \textit{send} category --- both the most
frequent intent in the dataset and the one with the highest numeral and
proper-noun density, consistent with the hypothesis that domain adaptation
disproportionately helps exactly the structurally hardest utterances.

\subsection{Acoustic Robustness}

Because financial voice interfaces are often used in mobile contexts, we synthetically evaluated our domain-adapted model against real-world acoustic degradations using the \textit{audiomentations} library. The original test set consists of clean read speech (42.58\% WER). When subjected to simulated GSM phone band-limiting (300Hz--3400Hz), the model exhibited negligible degradation, achieving a 46.70\% WER. Mild simulated room reverberation degraded performance to 51.37\%. Under additive background street noise, the model degraded gracefully to 54.40\% WER at 10dB SNR, though it unsurprisingly failed at 0dB SNR (103.85\% WER) where the signal and noise are of equal power. The minimal degradation under phone band-limiting confirms the model's suitability for mobile accessibility applications.

\subsection{Data Efficiency}

A critical question for low-resource languages is determining the minimum data volume required to achieve usable domain adaptation. We trained our LoRA configuration across increasing subset sizes of the training data ($N \in \{50, 100, 150, 200, 300, 403\}$). At merely 50 examples, Word Error Rate dropped from 131.0\% (zero-shot) to 73.9\%, and Transaction Success Rate (TSR) reached 18.3\%. As data increased, performance consistently improved, plateauing at $N=300$ with a peak TSR of 38.3\%, before slightly degrading at $N=403$ (33.3\% TSR, 43.1\% WER). This demonstrates that significant domain adaptation in highly specialized, constrained vocabularies (like financial digits) can be achieved with fewer than 300 transcribed utterances when utilizing parameter-efficient fine-tuning on large foundational models.

% ── 6. Error Analysis ─────────────────────────────────────────────────────────
\section{Error Analysis}

\subsection{Numeral Confusions}
While domain adaptation increased numeral exact-match accuracy from 0.0\% to 73.9\%, we analyzed the remaining errors to understand the failure modes of the adapted model. In 25 out of the 60 test utterances, the model transcribed exactly one numeral incorrectly. We observed three distinct patterns in these confusions:
\begin{enumerate}
    \item \textbf{Zero Insertion/Deletion:} The most common failure mode involves missing or adding trailing zeros (e.g., transcribing \dev{५०००} (5,000) as \dev{५००००} (50,000), or \dev{६५००००} (650,000) as \dev{६५०००} (65,000)). This indicates the model occasionally struggles to distinguish the acoustic duration of repeating zero words in Nepali (like \textit{hajaar} vs \textit{laakh}).
    \item \textbf{Prefix Hallucinations:} The model occasionally prepends the digit \dev{५} (5) to amounts (e.g., \dev{८००} transcribed as \dev{५८००}). This is likely an artifact of the acoustic similarity between the Nepali word for "Rs." (\textit{Rupaiya}) and the number 5 (\textit{Paanch}), which frequently co-occur in the training data.
    \item \textbf{Similar Sounding Digits:} We observed phonetic confusions such as \dev{१२} (12) and \dev{१} (1) (e.g., \dev{२००००} transcribed as \dev{१२००००}).
\end{enumerate}
These specific failures highlight why ASR systems for financial applications cannot rely on raw transcriptions alone and require secondary confirmation UI steps before executing transactions.

\subsection{Taxonomy Gaps}
Qualitative inspection of the ``other'' category (Table~\ref{tab:intent})
surfaces a genuine taxonomy gap rather than annotation noise. A recurring
pattern --- third-person statements describing funds already received,
e.g.\ \dev{अर्चना श्रेष्ठले कुमारी बैंकमा रु ४७५० प्राप्त गरेका छन्} (``Archana
Shrestha received Rs.\ 4750 at Kumari Bank'') --- does not map cleanly onto
the send/load/balance taxonomy used elsewhere in this work, since it is
neither a command nor strictly a balance query. We did not anticipate this
pattern during dataset design; it appears to have entered the corpus
because some contributors interpreted ``describe a financial transaction''
prompts more broadly than intended.

We treat this as a useful negative result: an intent classifier trained
naively on the present three-class scheme would silently misclassify this
``receive'' pattern, and any production deployment of the application's
NLP layer should add a fourth intent class or explicitly reject
out-of-taxonomy utterances rather than forcing a best-effort match. We did
not retrain the application's intent parser to add this class in the
current version; this is identified as future work in Section~7.

A second, smaller pattern in the ``other'' category is unrelated
demographic statements (e.g.\ stating one's age) appearing in place of an
expected financial command. The frequency is too low ($n=4$ of 60 test
utterances) to support a strong claim about cause, but it is consistent
with occasional prompt-following drift during data collection rather than
a model artifact, since it appears identically in both zero-shot and
domain-adapted transcriptions of the same audio.

% ── 7. System Architecture and Deployment ──────────────────────────────────────
\section{System Architecture and Deployment}

The application is a Next.js web app deployed on Vercel, communicating with
a Supabase-backed PostgreSQL database (authentication, wallet balances, and
an atomic SQL transfer function to prevent race conditions on concurrent
transactions). Voice commands are captured client-side, sent to a
self-hosted FastAPI inference server (Hugging Face Spaces, Docker runtime)
running the LoRA-adapted model, and the resulting transcript is parsed by a
two-stage rule-and-confidence intent classifier before a spoken
confirmation is presented to the user. Self-hosting the inference endpoint,
rather than depending solely on a third-party serverless inference API,
was adopted after observing unreliable cold-start behavior under the free
tier of hosted inference --- a practical deployment lesson as relevant as
the modeling result itself for a system intended for real accessibility use.

\subsection{Informal Usability Observations}

Prior to the current LoRA-adapted model and voice-interaction redesign
described in this report, an early prototype of the application was
informally tested with 3--5 sighted volunteers using blindfolds to
simulate visual impairment. No formal protocol, timing instrumentation,
or written consent process was used at this stage --- this was
exploratory, not a controlled study, and we report it as such rather
than overstating its rigor.

Results were mixed: testers were able to complete some voice-driven
tasks, but encountered friction at several points in the interaction
flow, consistent with the prototype's reliance on the zero-shot base
ASR model (Table~\ref{tab:results}) and a less structured confirmation
flow than the current version. We do not have preserved logs or timing
data from these sessions and do not report quantitative metrics from
them.

This informal round directly motivated three subsequent design changes
documented elsewhere in this report: (1) the move to a domain-adapted
ASR model rather than the zero-shot base, given the qualitative
incoherence observed in early transcriptions (Section~5); (2) the
explicit voice-confirmation step before any financial action is
executed, added after testers expressed uncertainty about whether a
command had been correctly understood before it was acted on; and (3)
the self-hosted inference deployment described above, motivated by
unreliable response latency observed during testing under the
previous hosted-inference setup.

A structured usability evaluation --- ideally with visually impaired
participants rather than blindfolded sighted volunteers, following a
protocol such as the System Usability Scale used by Ilyasa et
al.~\cite{ilyasa2023} --- on the current version of the system remains
important future work and is the most significant remaining gap
between this report's model-level contribution and a complete
accessibility evaluation of the deployed application.

% ── 8. Limitations and Future Directions ───────────────────────────────────────
\section{Limitations and Future Directions}

\textbf{Incomplete baseline comparison.} The general-domain Nepali
fine-tune baseline failed to load during the benchmark run reported here
(Table~\ref{tab:results}). The present comparison is therefore zero-shot
vs.\ domain-adapted only; a complete three-way comparison against a strong
general-domain fine-tune is necessary before claiming the gain is specific
to financial-domain adaptation rather than to fine-tuning on \textit{any}
Nepali data. This is the highest-priority follow-up experiment.

\textbf{Dataset scale.} 403 utterances (303 for training) is small relative
to standard ASR fine-tuning corpora. The results should be read as a
feasibility demonstration --- domain adaptation provides large relative
gains even at this scale --- not as a claim of state-of-the-art absolute
performance. NumAcc of 73.9\% indicates meaningful remaining error on
numeral transcription, the single most safety-critical metric for a
payment application; this is not yet production-ready accuracy.

\textbf{Single-annotator transcription.} Transcripts were manually
verified but not independently double-annotated; no inter-annotator
agreement statistic is available.

\textbf{Intent taxonomy gap.} As discussed in Section~6, a recognizable
``funds received'' utterance pattern is not represented in the current
three-class intent scheme. Future work should add this as an explicit
fourth class and re-evaluate the application's NLP intent parser
specifically (as distinct from the ASR transcription quality evaluated in
this report).

\textbf{No formal accessibility evaluation.} No usability study has yet
been conducted with visually impaired users, the system's intended
audience. A System Usability Scale (SUS) study, following the methodology
used by Ilyasa et al.~\cite{ilyasa2023}, is planned future work and would
substantially strengthen the application-level (as opposed to model-level)
contribution of this project.

\textbf{Training dynamics not logged.} Per-checkpoint validation WER/loss
during training was not preserved for this report; a re-run with logged
checkpoints would allow direct assessment of whether the model overfit
within the chosen training step budget.

% ── 9. Conclusion ────────────────────────────────────────────────────────────
\section{Conclusion}

We presented NepFinSpeech-403, the first domain-specific Nepali financial
speech dataset, and showed that LoRA fine-tuning of Whisper large-v2 on
this small corpus produces a large, statistically robust improvement over
the zero-shot baseline: a 67.5\% relative WER reduction, recovery of basic
output coherence in a domain where the base model otherwise produces
incoherent transcriptions, and improvement on 59 of 60 individual test
utterances ($p = 1.7\times10^{-12}$). The gain is consistent across command
types, and error analysis surfaced a genuine gap in our intent taxonomy
rather than merely confirming existing assumptions. The system is deployed
as a working voice-first eWallet application, demonstrating that the
distance between a research result and a usable accessibility tool is
bridgeable with modest additional engineering. The principal remaining gap
--- a complete comparison against a general-domain fine-tuned baseline ---
is identified as the immediate next experiment.

The full codebase, training pipeline, dataset, model weights, and this
report are available at
\url{https://github.com/subedibiraj/speakpay}.

\bibliographystyle{unsrt}
\begin{thebibliography}{99}

\bibitem{who2019}
World Health Organization,
``World report on vision,''
2019.

\bibitem{regmi2021}
S.~Regmi and B.~K.~Bal,
``An end-to-end speech recognition for the Nepali language,''
in \textit{Proc. 18th International Conference on Natural Language Processing}, 2021.

\bibitem{paudel2023}
S.~Paudel, B.~K.~Bal, and D.~Shrestha,
``Large vocabulary continuous speech recognition for Nepali language using CNN and Transformer,''
in \textit{Proc. 4th Conference on Language, Data and Knowledge (LDK)}, 2023.

\bibitem{radford2022}
A.~Radford, J.~W.~Kim, T.~Xu, G.~Brockman, C.~McLeavey, and I.~Sutskever,
``Robust speech recognition via large-scale weak supervision,''
in \textit{Proc. ICML}, 2023.

\bibitem{hu2022lora}
E.~J.~Hu, Y.~Shen, P.~Wallis, Z.~Allen-Zhu, Y.~Li, S.~Wang, L.~Wang, and W.~Chen,
``LoRA: Low-rank adaptation of large language models,''
in \textit{Proc. ICLR}, 2022.

\bibitem{ilyasa2023}
I.~Ilyasa, R.~Fauzi, and S.~Suakanto,
``Creating a mobile e-wallet for visual impairment using an inclusive design approach,''
in \textit{Proc. ICOEINS}, 2023.

\bibitem{zhou2023}
Z.~Zhou, T.~Sharma, L.~Emano, S.~Das, and Y.~Wang,
``Iterative design of an accessible crypto wallet for blind users,''
\textit{arXiv preprint arXiv:2306.06261}, 2023.

\end{thebibliography}

\end{document}
